\addtotoc{Abstract}
\begin{abstract}
	\addtocontents{toc}{}

	\hspace{1cm} Autonomous penetration testing agents powered by large language models have shown rapid progress in recent years, yet published evaluations consistently point to two recurring weaknesses: agents explore candidate vulnerabilities too narrowly, and they struggle to reason about the prerequisite relationships between different stages of an attack. These weaknesses limit how far current systems can operate reliably against realistic, multi-step web, API, and network targets without close human supervision.

	\vspace{0.3cm}
	\hspace{1cm} This report documents the first stage of an investigation into whether explicitly modeling the dependency relationships between attack steps --- rather than treating candidate vulnerabilities as an unordered list --- can improve how autonomous agents explore and rank their actions during a penetration test. As an initial step, the authors reviewed a body of recent literature on LLM-based offensive security agents, covering single-agent systems, multi-agent orchestration frameworks, planning-augmented approaches, and the benchmarks used to assess them. Two recurrent failure modes --- insufficient exploration breadth and weak prerequisite reasoning --- surfaced across independently developed systems. These observations motivate the research direction described in this report, which goes by the name RedGrid.

	\vspace{0.3cm}
	\hspace{1cm} At this stage, the work is preliminary: a structured review of the current state of autonomous penetration testing research and an early architectural direction motivated by the gaps observed in the literature. The remaining work over the coming months focuses on refining this direction, implementing and assessing it against established benchmarks, and determining whether dependency-aware exploration meaningfully improves autonomous exploitation performance.

\end{abstract}

\clearpage
\setstretch{1.5}
